% 第1章 序論
% 対応する草稿: research/thesis-drafts/chapter1-introduction-draft.md

\subsection{背景}\label{subsec:background}

% TODO: chapter1-introduction-draft.md §1.1 を LaTeX 化

\subsection{問題意識}\label{subsec:problem}

% TODO: chapter1-introduction-draft.md §1.2 を LaTeX 化

\subsection{理論的視点: 動的認識論理から見たLLMの内省}\label{subsec:del-motivation}

% TODO: chapter1-introduction-draft.md §1.2' を LaTeX 化
% DEL フレーミング（教員承認後に正式化）

\subsection{研究目的と問い}\label{subsec:rq}

% TODO: chapter1-introduction-draft.md §1.3 を LaTeX 化
% RQ0〜RQ4 を箇条書きで整理

\subsection{本論文の構成}\label{subsec:structure}

% TODO: 各章の概要を 1〜2 文ずつ
